\KOMAoption{paper}{a4, portrait}
\areaset{267mm}{180mm}
\fancyheadoffset{0pt} % обновляем хедер и футер

\newgeometry{
  top=14mm,
  bottom=11mm,
  right=10mm,
  left=22mm,
}

\begin{landscape}

%\setlength{\headsep}{7.5mm} 
%\setlength{\topmargin}{-27mm}
\setlength{\footskip}{15.5pt}

\color{uniblue}{\section[(рекомендуемое) Подключение аналоговых цепей трансформаторов тока и напряжения к устройству]{(рекомендуемое)\\Подключение аналоговых цепей трансформаторов тока и напряжения к устройству}\label{app:analogs}}
\color{black}

\begin{figure}[H]
  \centering
  \includegraphics[width=1.0\textwidth,height=1.0\textheight,keepaspectratio]{\apppath/Приложение. Аналоги/_latex/img/1.png}
  \captionof{figure}{Подключение цепей ТТ и ТН к аналоговому модулю устройству (в случае двухобмоточного трансформатора)}\label{appcard:img1}
 \end{figure}

 \begin{figure}[H]
  \centering
  \includegraphics[width=1.3\textwidth,height=1.3\textheight,keepaspectratio]{\apppath/Приложение. Аналоги/_latex/img/2.png}
  \captionof{figure}{Подключение цепей ТТ и ТН к аналоговому модулю устройству (в случае двухобмоточного трансформатора с расщепленной обмоткой НН)}\label{appcard:img2}
 \end{figure}

 \begin{figure}[H]
  \centering
  \includegraphics[width=1.3\textwidth,height=1.3\textheight,keepaspectratio]{\apppath/Приложение. Аналоги/_latex/img/3.png}
  \captionof{figure}{Подключение цепей ТТ и ТН к аналоговому модулю устройству (в случае трехобмоточного трансформатора)}\label{appcard:img3}
 \end{figure}

\clearpage
\end{landscape}
\restoregeometry


